\chapter{Behavioural and Fit}\label{ch:iv:behavioural-and-fit}

``Tell me about the worst loss you have been responsible for.'' The candidate who answers with a loss that was someone else's
fault, or with a loss that was not really a loss, has answered a different question. The one who names the number, what she
believed at the time, what she did in the first hour and what she changed afterwards has answered the one being scored.
\omterm{def:iv:behavioural-and-fit:question}{Behavioural questions} are the part of the loop candidates prepare least and interviewers score most consistently, because they
are usually structured: the same questions for every candidate, answers scored against anchors (One Quant Book~16, chapter~10).

\section{What behavioural questions measure}

\begin{definition}[Behavioural question]\label{def:iv:behavioural-and-fit:question}
A \emph{behavioural question}\index{behavioural question} asks the candidate to describe what they did in a past situation of a
given kind (``tell me about a time when\dots''), on the premise that past behaviour in similar situations predicts future
behaviour; a situational question asks instead what the candidate would do in a described situation.
\end{definition}

The premise has evidence behind it. Structured interviews built from past-behaviour questions were compared with unstructured
ones as early as the patterned behaviour description interview (Janz, 1982), and situational questions were developed at the same
time (Latham, Saari, Pursell and Campion, 1980). A meta-analysis found both formats valid, with past-behaviour questions scored on
descriptively anchored scales reaching a higher mean validity than situational ones (0.63 against 0.47; Taylor and Small, 2002).
What the interviewer records is therefore specific: what the candidate did, not what the candidate thinks one should do.

\section{Structuring an answer}

\begin{definition}[STAR structure]\label{def:iv:behavioural-and-fit:star}
The \emph{STAR structure}\index{STAR structure} orders an answer to a \omterm{def:iv:behavioural-and-fit:question}{behavioural question} as the situation (the context, in two
sentences), the task (what the candidate was responsible for), the action (what the candidate did, in the first person, with the
reasoning), and the result (what happened, with a number, and what the candidate learned or changed).
\end{definition}

\begin{method}[Preparing behavioural answers]\label{met:iv:behavioural-and-fit:prepare}
\begin{enumerate}
\item Choose six stories from your own work or study: a success, a failure or loss, a disagreement, a mistake you caught, a time you
led, and a time you changed your mind.
\item Write each in the \omterm{def:iv:behavioural-and-fit:star}{STAR structure}, with the action taking most of the time and a number in the result.
\item Check each story for the first person: ``we decided'' tells the interviewer nothing about you.
\item For the failure and the loss, separate the quality of the decision from the outcome (One Quant Book~16, chapter~26, on
outcome bias), and say what you changed.
\item Practise aloud to two minutes each; stop when the result is told.
\end{enumerate}
\end{method}

\section{Why trading, why this firm, why this role}

These questions are scored for specificity and consistency with the rest of the loop. A good answer names what the candidate
likes about the work itself (fast feedback, decisions under uncertainty, building systems that are measured by money) with an
example from the candidate's own experience, and why this kind of firm (a market maker's graduate class, a systematic fund's
research culture, a bank's client flow) fits it. ``I want to make money'' is honest and incomplete; ``I like puzzles'' says
nothing about trading. The firms' own descriptions of the roles are in One Quant Book~17, chapters 16 to~25.

\section{The loss question and integrity scenarios}

The loss question tests three things: whether the candidate owns the loss, whether the candidate can tell a bad decision from a
bad outcome, and what the candidate did next. A candidate without trading experience answers with the nearest equivalent (a
competition, a project that failed, a bet in a game) and the same structure.

\begin{definition}[Integrity scenario]\label{def:iv:behavioural-and-fit:integrity}
An \emph{integrity scenario}\index{integrity scenario} is an interview question that describes a situation with a conflict between
a candidate's interest and a rule or duty (possible inside information, an error that only the candidate has seen, pressure to
present a result one doubts) and asks what the candidate would do.
\end{definition}

\omterm{def:iv:behavioural-and-fit:integrity}{Integrity scenarios} have right answers, and the right answers are unglamorous: stop, do not trade, tell the person whose job it is
to know (a manager, compliance, risk), and write it down. They engage the rules of One Quant Book~16, chapter~16: insider dealing
and the unlawful disclosure of inside information are prohibited (the EU Market Abuse Regulation, articles 8, 10 and~14, and their
equivalents elsewhere), personal trading goes through pre-clearance and the restricted and watch lists, and errors are reported at
once. An interview is not legal advice and does not need legal detail; it needs the reflex.

The arithmetic behind a loss question is the same as in \cref{iq:iv:behavioural-and-fit:5}: put the loss in units of the book's
standard deviation and ask how often such a day should come. \Cref{fig:iv:behavioural-and-fit:tails} shows why the answer
depends on the tails far more than on the standard deviation.

\begin{omfigure}
\begin{tikzpicture}
\begin{semilogyaxis}[width=9.5cm,height=5.4cm,xlabel={\small size of the daily loss, $k$ standard deviations},
  ylabel={\small trading days between such losses},xmin=1,xmax=5,ymin=1,ymax=1e7,
  tick label style={font=\footnotesize},grid=major,grid style={gray!20},table/col sep=comma,
  legend cell align=left,legend pos=north west,legend style={font=\scriptsize}]
\addplot[omDef,thick] table[x=k,y=normal]{figdata/interviews/28-behavioural-and-fit/tails.csv};
\addplot[omThm,thick,dashed] table[x=k,y=t3]{figdata/interviews/28-behavioural-and-fit/tails.csv};
\addplot[omProp,thick,densely dotted] table[x=k,y=t5]{figdata/interviews/28-behavioural-and-fit/tails.csv};
\legend{normal,Student $t$ (3 d.f.),Student $t$ (5 d.f.)}
\end{semilogyaxis}
\end{tikzpicture}

\omcaption{Average number of trading days between daily losses of at least $k$ standard deviations, when daily P\&L is normal
or follows a Student $t$ distribution scaled to the same standard deviation. A three-standard-deviation loss comes once in
741 days under the normal and once in 144 under the $t$ with three degrees of freedom; at five standard deviations the gap is a
factor of several thousand. Data: \texttt{fig\_iv\_tails.py}.}\label{fig:iv:behavioural-and-fit:tails}
\end{omfigure}

\section{Worked answers}

\begin{example}[A mistake you caught, in the STAR structure]\label{ex:iv:behavioural-and-fit:caught}
\emph{``Tell me about a time you found an error in your own work after you had shared it.''}

\emph{Situation.} ``In my master's project I sent my supervisor a backtest of a momentum signal on thirty commodity futures
with a Sharpe ratio of 1.3.'' \emph{Task.} ``I was responsible for the data pipeline, including the roll from one contract
to the next.'' \emph{Action.} ``Two days later, plotting the continuous price series for a different reason, I saw jumps on
roll dates: I had joined the contracts without adjusting for the gap between them, so the signal was partly trading the
roll. I wrote to my supervisor that afternoon, before rerunning anything, to say the number was wrong and why. Then I
rebuilt the series with ratio-adjusted rolls, wrote a test that fails if any daily return on a roll date exceeds five
standard deviations, and reran.'' \emph{Result.} ``The Sharpe ratio fell to 0.6. The project's conclusion changed from
`strong' to `modest and cost-sensitive', and the roll test is now the first thing I write on any futures data.''

The answer takes about ninety seconds, is in the first person throughout, discloses before fixing, gives both numbers, and
ends with a change of practice (One Quant Book~7, chapter~2, on continuous futures series).
\end{example}

\begin{example}[Why this firm, weak and strong]\label{ex:iv:behavioural-and-fit:why}
A software engineer applying to an options market maker. \emph{Weak:} ``You are a leading firm with great technology and
smart people, and I want to work on challenging problems.'' Every word could be said to any firm, and the interviewer learns
nothing. \emph{Strong:} ``I've spent two years on a payments system where the hard problem was correctness under load. At a
market maker the same problem has a clock on it: a quote that is correct but late is wrong. I want to work where the latency
of my code is measured in money, and your engineers' public talks on measuring tail latency are the kind of work I want to
do.'' It names the candidate's experience, a property of the business that the candidate understands, and a specific,
checkable reason for this firm, without flattery.
\end{example}

\section{Question bank}

\begin{interviewq}[$\star$ \iqroles{trader, researcher} \iqfirm{market maker}]\label{iq:iv:behavioural-and-fit:1}
Why trading? Give the answer you would give, and say what makes it score.
\end{interviewq}

\begin{interviewq}[$\star$ \iqroles{researcher, developer} \iqfirm{systematic fund}]\label{iq:iv:behavioural-and-fit:2}
Why a systematic fund rather than a bank or a market maker?
\end{interviewq}

\begin{interviewq}[$\star$ \iqroles{developer, mle} \iqfirm{any}]\label{iq:iv:behavioural-and-fit:3}
Tell me about a time you disagreed with a teammate about a technical decision.
\end{interviewq}

\begin{interviewq}[$\star\star$ \iqroles{trader, risk} \iqfirm{multi-manager fund}]\label{iq:iv:behavioural-and-fit:4}
Tell me about the worst loss you have been responsible for. (Answer it for a candidate with two years on a desk, and for a
graduate.)
\end{interviewq}

\begin{interviewq}[$\star\star$ \iqroles{trader, risk} \iqfirm{bank}]\label{iq:iv:behavioural-and-fit:5}
Your book lost 2.4 million yesterday; its daily P\&L has a standard deviation of 0.8 million. How often should a day like that
happen, and what do you tell your head of desk this morning?
\end{interviewq}

\begin{interviewq}[$\star\star$ \iqroles{developer} \iqfirm{proprietary firm}]\label{iq:iv:behavioural-and-fit:6}
Tell me about a bug of yours that reached production.
\end{interviewq}

\begin{interviewq}[$\star\star$ \iqroles{trader, researcher, developer} \iqfirm{any}]\label{iq:iv:behavioural-and-fit:7}
``Do you have any questions for me?'' What do you ask a trader, a researcher and a head of desk?
\end{interviewq}

\begin{interviewq}[$\star\star\star$ \iqroles{trader} \iqfirm{asset manager}]\label{iq:iv:behavioural-and-fit:8}
A friend who works at a listed company tells you over dinner that its results, due next week, are much better than expected. What do
you do, and what must you avoid?
\end{interviewq}

\begin{interviewq}[$\star\star\star$ \iqroles{researcher, mle} \iqfirm{systematic fund}]\label{iq:iv:behavioural-and-fit:9}
The night before an investment committee, you find that a colleague's backtest, which the committee will rely on, includes a
look-ahead error. The colleague is senior and unreachable. What do you do?
\end{interviewq}

\begin{interviewq}[$\star\star\star$ \iqroles{trader, developer} \iqfirm{market maker}]\label{iq:iv:behavioural-and-fit:10}
You notice that an order you sent an hour ago had an extra zero in its size; it filled, and the position has since made money.
Nobody else has noticed. What do you do?
\end{interviewq}

\begin{omsources}
\item T.~Janz, ``Initial comparisons of patterned behavior description interviews versus unstructured interviews'', \emph{Journal
of Applied Psychology} 67(5), 1982, 577--580.
\item G.~P.~Latham, L.~M.~Saari, E.~D.~Pursell and M.~A.~Campion, ``The situational interview'', \emph{Journal of Applied
Psychology} 65(4), 1980, 422--427.
\item L.~Taylor and B.~Small, ``Asking applicants what they would do versus what they did do: a meta-analytic comparison of
situational and past behaviour employment interview questions'', \emph{Journal of Occupational and Organizational Psychology}
75(3), 2002, 277--294.
\item Regulation (EU) No 596/2014 (Market Abuse Regulation), articles 7, 8, 10 and~14; One Quant Book~16, chapters 10, 16 and~26.
\end{omsources}
